%% =====================================================================
%% CANONICAL TEMPLATE — IEEE conference (IEEEtran, two-column)
%% Copy this file, keep the structure, swap the content.
%% Build:  writter-ai/skills/academic-paper/scripts/build.sh example-ieee.tex
%% Class docs: texdoc IEEEtran  ·  https://www.ieee.org/conferences/publishing/templates
%% =====================================================================
\documentclass[conference]{IEEEtran}
% For a journal (single/double column, per-journal) use e.g.:
%   \documentclass[journal]{IEEEtran}
% Camera-ready page limits: strip \IEEEoverridecommandlockouts only if needed.

\usepackage{cite}                 % IEEE-style [1],[2]-[4] citation compression
\usepackage{amsmath,amssymb,amsfonts}
\usepackage{algorithmic}
\usepackage{graphicx}
\usepackage{textcomp}
\usepackage{xcolor}
\usepackage{booktabs}             % professional tables (\toprule \midrule \bottomrule)
% hyperref LAST; IEEEtran-friendly options:
\usepackage[hidelinks,breaklinks=true]{hyperref}

\begin{document}

\title{A Concise, Specific Paper Title\\
{\footnotesize \textsuperscript{}Optional subtitle or venue note}}

% ---- Authors (IEEEtran affiliation blocks) ----
\author{
\IEEEauthorblockN{First A. Author}
\IEEEauthorblockA{\textit{Dept. Name} \\
\textit{Organization}\\
City, Country \\
first.author@org.edu}
\and
\IEEEauthorblockN{Second B. Author}
\IEEEauthorblockA{\textit{Dept. Name} \\
\textit{Organization}\\
City, Country \\
second.author@org.edu}
}

\maketitle

\begin{abstract}
State the problem in one sentence, the gap in the second, your approach in the
third, and the headline quantitative result in the fourth. An abstract is
150--250 words, contains no citations, and every claim here is proven in the body.
Replace this paragraph with yours.
\end{abstract}

\begin{IEEEkeywords}
keyword one, keyword two, keyword three, keyword four
\end{IEEEkeywords}

\section{Introduction}
Open with the problem and why it matters, then the specific gap in prior work,
then your contribution. Close the introduction with an explicit contribution
list, which reviewers look for:
\begin{itemize}
\item A clearly stated method/algorithm (Section~\ref{sec:method}).
\item An empirical result: e.g., we improve X by \emph{Y\%} over baseline~\cite{ref_baseline}.
\item An analysis/ablation isolating the source of the gain (Section~\ref{sec:results}).
\end{itemize}

\section{Related Work}
Group prior art into 2--3 threads; for each, state what it does and its limitation
relative to your setting~\cite{ref_survey,ref_method}. Do not merely list papers.

\section{Method}\label{sec:method}
Describe the approach precisely enough to reproduce. Define notation once. A
displayed equation is numbered and referenced as Eq.~\eqref{eq:objective}:
\begin{equation}\label{eq:objective}
\mathcal{L}(\theta) \;=\; \frac{1}{N}\sum_{i=1}^{N} \ell\!\left(f_\theta(x_i), y_i\right) + \lambda\,\lVert\theta\rVert_2^2 .
\end{equation}
State assumptions explicitly. If there is an algorithm, present it in an
\texttt{algorithmic} block and reference its line numbers in the text.

\section{Experiments}\label{sec:exp}
Specify datasets, baselines, metrics, and the exact protocol (splits, seeds,
hardware). Every number in the paper must be reproducible from this description.

\section{Results}\label{sec:results}
Lead with the headline table, then analysis. Table~\ref{tab:main} spans one column;
for a full-width float use \verb|table*| and \verb|figure*|.
\begin{table}[t]
\caption{Main results. Best per column in \textbf{bold}. Report mean$\pm$std over seeds.}
\label{tab:main}
\centering
\begin{tabular}{lcc}
\toprule
Method & Metric A ($\uparrow$) & Metric B ($\downarrow$) \\
\midrule
Baseline~\cite{ref_baseline} & 71.2 & 0.184 \\
Ours                         & \textbf{78.6} & \textbf{0.142} \\
\bottomrule
\end{tabular}
\end{table}

% Full-width figure example (spans both columns, top of page):
%\begin{figure*}[t]
%  \centering
%  \includegraphics[width=\textwidth]{figure.pdf}
%  \caption{Caption. Explain what the reader should conclude, not just what is shown.}
%  \label{fig:overview}
%\end{figure*}

\section{Conclusion}
Restate the contribution and the single most important number. State one concrete
limitation and one next step. No new claims here.

\section*{Acknowledgment}
Optional. Funding, compute, and people who helped but are not authors.

% ---- References ----
% IEEEtran.bst is fetched automatically by tectonic. Put entries in refs.bib.
\bibliographystyle{IEEEtran}
\bibliography{refs}

\end{document}
